\section{Kernel-Level Analysis: How Quantum Is the Kernel Advantage?}\label{sec:kernel}
The audit's most robust positive was the quantum-kernel SVM out-ranking its random-feature surrogate
(Section~\ref{sec:audit-results}), the one comparison to survive both FDR and family-wise Holm
correction. That control, however, differed from the QSVM in four ways besides the feature map: it was
trained on 1{,}200 rather than 2{,}000 rows, it approximated its kernel with 512 Nystr\"om components,
it used a linear rather than a kernel SVM, and its bandwidth $\gamma=1$ was applied untuned to raw
features on a different scale. This section removes each difference in turn, examines how the result
depends on hyperparameter selection, tests it on a second official split, and asks whether what remains
requires the quantum state. The diagnostics are those of Section~\ref{sec:kdiag}.

\subsection{The surrogate gap is set by the classical bandwidth search}\label{sec:kgap}
Table~\ref{tab:fair} places every kernel on the identical 2{,}000 NSL-KDD training rows, with exact Gram
matrices and one kernel SVM solver. Removing the first three confounds, with the bandwidth still fixed at
$\gamma=1$, narrows the published ROC-AUC gap from $0.127$ to $0.085$. Hyperparameter selection then
decides the outcome. With the bandwidth searched over $\gamma\le3$, cross-validation picks $\gamma=0.3$ for
the raw-feature RBF kernel, whose test ROC-AUC falls to $0.783$: an apparent quantum advantage of $0.167$
(95\% CI $[0.161,0.173]$). With the same cross-validation over $\gamma\le30$ it picks $\gamma=30$, and the RBF
kernel reaches $0.940$. The quantum kernel then leads by only $0.007$ in ROC-AUC ($[0.006,0.009]$) and
\emph{trails} at the operating point, where the classical kernel detects $0.455$ of attacks at $1\%$ false
positives against $0.442$ ($\Delta=-0.013$, $[-0.020,-0.005]$, $p=0.002$). The advantage the audit reported
is therefore not a property of the kernels. It appears or disappears with the width of the classical
baseline's search grid, a special case of the bandwidth sensitivity documented for quantum kernels by
Shaydulin and Wild~\cite{shaydulin2022bandwidth}, here acting on the \emph{classical} side of the
comparison.

@@TAB_FAIR@@

\subsection{Cross-validation cannot see the novel-attack shift}\label{sec:kcv}
Figure~\ref{fig:landscape} shows why the grid matters so much. Within the training distribution every
kernel scores between $0.95$ and $0.99$ at every bandwidth for $C=1$, so the cross-validation surface is nearly flat
and its maximum is close to arbitrary. On the official test set, whose 17 attack types never occur in
training, the raw-feature RBF kernel is U-shaped in $\gamma$: strong at very wide and very narrow
bandwidths (up to $0.945$), and weakest at intermediate ones ($0.740$ at worst across the grid), which is exactly where a
grid capped at $\gamma=3$ leads cross-validation. Across the 40 settings of the grid the RBF kernel's
standard CV score is uncorrelated with its test ROC-AUC ($r=-0.16$).

This motivates \emph{shift-aware model selection}: leave-attack-types-out cross-validation, in which each
attack type lives in exactly one fold, so every validation fold contains attack types its training fold
never saw. It reproduces the condition of the official split using training data alone. Its score
correlates with test ROC-AUC at $r=+0.79$ for the RBF kernel, and it selects robust settings for every
kernel family, with test ROC-AUC between $0.939$ and $0.947$ (Table~\ref{tab:fair}). Under this rule the
quantum projected kernel also reaches the highest operating-point detection we observe,
$\tpr$@$1\%\fpr=0.499$, above the RBF kernel's $0.466$ ($\Delta=+0.033$, $[0.021,0.067]$, $p<0.001$). We do
not read this as a quantum advantage, because its sign depends on the selection rule: under standard CV
the same comparison favours the classical kernel by $0.013$. Having now evaluated three selection rules,
reporting only the one on which quantum wins would be precisely the researcher degree of freedom this
paper warns against.

@@FIG_LANDSCAPE@@

\subsection{A second official split}\label{sec:kunsw}
If the collapse of standard CV is caused by novel attack types, it should vanish on a dataset whose test
set has none. UNSW-NB15 is the only other benchmark here with an official split, and every one of its
nine attack categories occurs in both training and test. We repeat the protocol on a seeded, stratified
sample of 2{,}000 training rows and 20{,}000 test rows; the training file is sorted by class, so its first
rows are all benign and cannot be used. Table~\ref{tab:twosplits} confirms the prediction: on UNSW-NB15
standard CV tracks test performance for every kernel ($r=+0.82$ to $+0.83$). The result for the quantum
kernel, however, reverses. Selected by the same rule, it is significantly worse than both classical kernels
on every metric: the classical phase kernel of Section~\ref{sec:kgeom} reaches test ROC-AUC $0.956$ against
$0.930$ ($\Delta=-0.026$, $[-0.028,-0.024]$) and detects $0.751$ of attacks at $1\%$ false positives against
$0.457$ ($\Delta=-0.294$, $[-0.318,-0.263]$), and even the raw-feature RBF kernel is ahead ($\Delta\tpr=-0.057$,
$[-0.087,-0.022]$, $p=0.005$). The one property that had favoured the quantum kernel on NSL-KDD, fewer poor
hyperparameter settings (worst test ROC-AUC $0.893$ against $0.740$), does not survive either: on
UNSW-NB15 its worst setting is the lowest of the three ($0.697$ against $0.787$ and $0.789$). We therefore do
not claim it as a property of quantum kernels. Shift-aware CV is not informative on UNSW-NB15, where the
dominance of a single attack category leaves only one usable fold, and we report no result from it there.

@@TAB_TWOSPLITS@@

\subsection{Geometry: the IQP kernel is close to its own classical phases}\label{sec:kgeom}
Table~\ref{tab:geom} gives the geometric difference on all four datasets, which requires no labels and no
hyperparameter selection. The angle-encoded fidelity kernel is reproduced by an RBF kernel to $g=1.06$ to
$1.07$ on every dataset, so by the bound of Huang et al.~\cite{huang2021power} it can hold no prediction
advantage over that classical kernel for any labelling. This is expected: under angle encoding every qubit
remains in a product state, and the kernel factorises into a product of cosines. The IQP kernels, which
use entangling ZZ gates, are genuinely far from standard kernels: against every RBF bandwidth and against
plain trigonometric features, $g$ ranges from $4.9$ to $12.2$. The same IQP kernels, however, lie within
$g=1.4$ to $2.4$ of a purely classical kernel built on the phases the ZZ feature map
writes~\cite{havlicek2019supervised}, on all four datasets, and their label-dependent model complexities
agree: on NSL-KDD the quantum projected kernel has $s_K=235$ and the classical phase kernel $s_K=236$, against
$481$ for the best raw-feature RBF. The empirical results are consistent with this geometry. On NSL-KDD the
phase kernel matches the quantum kernel at the operating point under a narrow grid ($0.446$ against
$0.444$, $p=0.89$) and under standard CV over the full grid ($0.441$ against $0.442$, $p=0.90$), and on
UNSW-NB15 it outperforms it. Whatever the IQP kernel contributes lives in the phases of its circuit, not
in the quantum state, and a classical kernel that uses those phases obtains it at a small fraction of the
cost.

@@TAB_GEOM@@

\subsection{Concentration and finite shots}\label{sec:kshots}
Kernel values concentrate as width grows, but mildly at these sizes. From 2 to 12 qubits the mean
off-diagonal IQP fidelity falls from $0.32$ to $0.047$, yet its variance falls only from $0.099$ to $0.032$
and flattens between 8 and 10 qubits, far from the exponential regime of Thanasilp
et al.~\cite{thanasilp2024concentration}. The kernels' test ROC-AUC peaks at 8 to 10 qubits ($0.944$ to
$0.946$ for the projected kernel) and dips only to $0.927$ at 12. Concentration therefore does not explain
the decline of the variational hybrid with width (Figure~\ref{fig:ablation}), which points instead to the
trainability and capacity of the variational model.

The projected kernel is robust to measurement noise (Table~\ref{tab:shots}). With 16 shots in each of three
measurement bases, 48 shots per sample in total, test ROC-AUC falls from $0.945$ to $0.941$, and AUPRC and
$\tpr$@$1\%\fpr$ do not move beyond their run-to-run spread; by 64 shots the result is indistinguishable
from the exact kernel. This follows from the construction: the projected kernel needs only single-qubit
expectations, each estimated with variance at most $1/M$, which the RBF then smooths, whereas a fidelity
kernel must resolve small global overlaps. The study models sampling noise only; gate and readout error on
the kernel circuit remain untested.

@@TAB_SHOTS@@

\subsection{Summary}\label{sec:ksum}
Once kernels are compared on identical data with equal and adequate tuning, no quantum-kernel advantage
survives the choice of selection rule or of dataset. The audit's strongest positive was produced by its
own classical surrogate: first by four confounds, then by a bandwidth grid too narrow for a distribution
shift that in-distribution cross-validation cannot see. Two findings of general use come out of the
exercise. Leave-attack-types-out cross-validation restores reliable model selection under novel-attack
shift, and the geometric difference identifies, without labels, when a quantum kernel is classically
reproducible, which here it is from its own circuit phases on every dataset.
